\documentclass[11pt,a4paper]{article}
\usepackage{amsmath,amssymb}
\usepackage{listings}
\usepackage{hyperref}
\usepackage[margin=1in]{geometry}
\usepackage{fancyhdr}

\pagestyle{fancy}
\fancyhf{}
\fancyhead[L]{Module 01: Getting Started with C}
\fancyhead[R]{C Programming Course}
\fancyfoot[C]{\thepage}

\lstset{
  language=C,
  basicstyle=\ttfamily\small,
  keywordstyle=\bfseries,
  commentstyle=\itshape,
  numbers=left,
  numberstyle=\tiny,
  frame=single,
  breaklines=true,
  tabsize=4
}

\title{Module 01: Getting Started with C}
\author{C Programming Course}
\date{}

\begin{document}
\maketitle
\thispagestyle{fancy}

\section{Introduction to C Programming}
C is a general-purpose, procedural programming language developed by Dennis Ritchie
at Bell Labs in 1972. It provides low-level access to memory and maps efficiently to
machine instructions, making it ideal for systems programming.

\section{Setting Up Your Development Environment}
You need a C compiler (GCC or Clang) and a text editor. On Linux:
\begin{verbatim}
sudo apt-get install gcc
\end{verbatim}
Verify the installation with \texttt{gcc --version}.

\section{Your First C Program}
\begin{lstlisting}
#include <stdio.h>

int main(void) {
    printf("Hello, World!\n");
    return 0;
}
\end{lstlisting}

\subsection{Key Elements}
\begin{itemize}
  \item \texttt{\#include <stdio.h>} -- includes the standard I/O header.
  \item \texttt{int main(void)} -- the program entry point.
  \item \texttt{printf} -- prints formatted output to the console.
  \item \texttt{return 0} -- indicates successful execution.
\end{itemize}

\section{Understanding the Compilation Process}
The compilation pipeline consists of four stages:
\begin{enumerate}
  \item \textbf{Preprocessing} -- expands macros and includes headers (\texttt{gcc -E}).
  \item \textbf{Compilation} -- translates C to assembly (\texttt{gcc -S}).
  \item \textbf{Assembly} -- converts assembly to object code (\texttt{gcc -c}).
  \item \textbf{Linking} -- combines object files into an executable.
\end{enumerate}

\begin{lstlisting}
/* Compile and run in one step */
// gcc hello.c -o hello && ./hello
\end{lstlisting}

\section{Basic Input/Output}
\begin{lstlisting}
#include <stdio.h>

int main(void) {
    int age;
    printf("Enter your age: ");
    scanf("%d", &age);
    printf("You are %d years old.\n", age);
    return 0;
}
\end{lstlisting}

Common format specifiers: \texttt{\%d} (int), \texttt{\%f} (float),
\texttt{\%c} (char), \texttt{\%s} (string).

\end{document}
